\section{Konfiguracje eksperymentalne}

W celu optymalizacji struktury sieci oraz algorytmu uczenia przygotowano
zestaw eksperymentów obejmujący różne architektury MLP oraz różne wartości
współczynnika uczenia.

\subsection{Analizowane wartości współczynnika uczenia}

W eksperymentach porównywano trzy wartości parametru \textit{learning rate}:
\begin{itemize}
    \item \(\eta = 0.001\),
    \item \(\eta = 0.01\),
    \item \(\eta = 0.1\).
\end{itemize}

\subsection{Analizowane struktury sieci}

Przygotowano kilka wariantów architektury sieci MLP różniących się liczbą
warstw ukrytych oraz liczbą neuronów w każdej z nich.

\begin{table}[H]
    \centering
    \begin{tabular}{lll}
\toprule
Nazwa konfiguracji & Struktura sieci & Uwagi \\
\midrule
small\_1\_hidden\_4 & \(13 \rightarrow 4 \rightarrow 3\) & wariant prosty \\
small\_1\_hidden\_8 & \(13 \rightarrow 8 \rightarrow 3\) & wariant prosty \\
medium\_1\_hidden\_13 & \(13 \rightarrow 13 \rightarrow 3\) & wariant bazowy \\
medium\_2\_hidden\_13\_8 & \(13 \rightarrow 13 \rightarrow 8 \rightarrow 3\) & wariant dwuwarstwowy \\
large\_2\_hidden\_32\_16 & \(13 \rightarrow 32 \rightarrow 16 \rightarrow 3\) & wariant większy \\
large\_3\_hidden\_64\_32\_16 & \(13 \rightarrow 64 \rightarrow 32 \rightarrow 16 \rightarrow 3\) & wariant największy \\
\bottomrule
\end{tabular}

    \caption{Konfiguracje architektur analizowane w etapie trzecim}
    \label{tab:network_configurations}
\end{table}

\subsection{Liczba kombinacji eksperymentalnych}

Każda architektura została połączona z każdą analizowaną wartością
współczynnika uczenia. Dzięki temu możliwe jest bezpośrednie porównanie
wpływu obu tych czynników na przebieg uczenia i końcową jakość klasyfikacji.
